% GENERATED FILE. Edit canonical lessons, then run npm run generate:books.
% canonical-source-hash: fnv1a64:5a0f6394e29d75e3
% canonical-lessons: HI-C427-udhar, HI-C427-qist, HI-C427-hisab, HI-C427-dalal, HI-C427-cukana

\chapter{Credit and Instalments}
\label{ch:hi-credit-and-instalments}
\hlchaptermodality{427}

\begin{chapteropening}
\textbf{By the end of this chapter:} \emph{I can buy on credit, pay in instalments, keep accounts and pay off a debt.}

Complete the last lesson of chapter 427: I can buy on credit, pay in instalments, keep accounts and pay off a debt.
\end{chapteropening}

\section[\texorpdfstring{\hi{उधार} (udhār)}{उधार (udhār)}]{\hi{उधार} (udhār) — credit; a loan}

\label{lesson:HI-C427-udhar}



\begin{quote}
Before the new one: say the Hindi for a balance, then the Hindi for a cash machine.
\end{quote}

\subsection*{You'll want to know: \hi{उधार}}
\textbf{\hi{उधार}} — \emph{udhār} — \textquotedblleft{}credit; a loan\textquotedblright{}.

\begin{quote}
\hi{दुकानदार} \hi{ने} \hi{मुझे} \hi{उधार} \hi{पर} \hi{चीनी} \hi{दी।}
\end{quote}

\begin{quote}
\emph{dukāndār ne mujhe udhār par cīnī dī.}
\end{quote}

\begin{quote}
The shopkeeper gave me sugar on credit.
\end{quote}

Masculine.

\begin{cousinweb}
\textbf{\hi{उधार}} (\emph{udhār}) is a tadbhava of Sanskrit \emph{uddhāra}, \textquotedblleft{}raising\textquotedblright{} and also \textquotedblleft{}a loan\textquotedblright{} (CDIAL 2018), which lists Hindi \emph{udhār}, \textquotedblleft{}a debt\textquotedblright{}.
\end{cousinweb}

\begin{grammarlens}[title={What you've built}]
One more word for getting around.
\end{grammarlens}

\subsection*{Your turn}
\begin{itemize}
\raggedright
  \item \emph{Say it:} \emph{udhār}
  \item \emph{Say it:} \emph{udhār}, once more
  \item \emph{Say it:} \emph{dukāndār ne mujhe udhār par cīnī dī.} — The shopkeeper gave me sugar on credit.
  \item \emph{Recall:} say the Hindi for a balance, then the Hindi for a cash machine, then say \emph{udhār} again
\end{itemize}

\begin{tcolorbox}[breakable,colback=teal!4,colframe=teal!35!black,title={Before you move on}]
What does \hi{उधार} mean? (\textquotedblleft{}Credit; a loan\textquotedblright{}.) Say it once more.
\end{tcolorbox}
\section[\texorpdfstring{\hi{क़िस्त} (qist)}{क़िस्त (qist)}]{\hi{क़िस्त} (qist) — an instalment}

\label{lesson:HI-C427-qist}



\begin{quote}
Before the new one: say the Hindi for a cash machine, then the Hindi for credit.
\end{quote}

\subsection*{You'll want to know: \hi{क़िस्त}}
\textbf{\hi{क़िस्त}} — \emph{qist} — \textquotedblleft{}an instalment\textquotedblright{}.

\begin{quote}
\hi{हम} \hi{हर} \hi{हफ़्ते} \hi{एक} \hi{क़िस्त} \hi{देते} \hi{हैं।}
\end{quote}

\begin{quote}
\emph{ham har hafte ek qist dete haiṁ.}
\end{quote}

\begin{quote}
We pay one instalment every week.
\end{quote}

Feminine.

\begin{cousinweb}
\textbf{\hi{क़िस्त}} (\emph{qist}) is usually traced to Arabic \emph{qisṭ}, \textquotedblleft{}a portion, a share\textquotedblright{}, through Persian: a debt paid in portions.
\end{cousinweb}

\begin{grammarlens}[title={What you've built}]
One more word for getting around.
\end{grammarlens}

\subsection*{Your turn}
\begin{itemize}
\raggedright
  \item \emph{Say it:} \emph{qist}
  \item \emph{Say it:} \emph{qist}, once more
  \item \emph{Say it:} \emph{ham har hafte ek qist dete haiṁ.} — We pay one instalment every week.
  \item \emph{Recall:} say the Hindi for a cash machine, then the Hindi for credit, then say \emph{qist} again
\end{itemize}

\begin{tcolorbox}[breakable,colback=teal!4,colframe=teal!35!black,title={Before you move on}]
What does \hi{क़िस्त} mean? (\textquotedblleft{}An instalment\textquotedblright{}.) Say it once more.
\end{tcolorbox}
\section[\texorpdfstring{\hi{हिसाब} (hisāb)}{हिसाब (hisāb)}]{\hi{हिसाब} (hisāb) — an account; a reckoning}

\label{lesson:HI-C427-hisab}



\begin{quote}
Before the new one: say the Hindi for credit, then the Hindi for an instalment.
\end{quote}

\subsection*{You'll want to know: \hi{हिसाब}}
\textbf{\hi{हिसाब}} — \emph{hisāb} — \textquotedblleft{}an account; a reckoning\textquotedblright{}.

\begin{quote}
\hi{दुकानदार} \hi{सब} \hi{ग्राहकों} \hi{का} \hi{हिसाब} \hi{रखता} \hi{है।}
\end{quote}

\begin{quote}
\emph{dukāndār sab grāhakoṁ kā hisāb rakhtā hai.}
\end{quote}

\begin{quote}
The shopkeeper keeps an account for every customer.
\end{quote}

Masculine.

\begin{cousinweb}
\textbf{\hi{हिसाब}} (\emph{hisāb}) is usually traced to Arabic \emph{ḥisāb}, \textquotedblleft{}counting, reckoning\textquotedblright{}, through Persian. \hi{हिसाब} \hi{करना} is to settle up.
\end{cousinweb}

\begin{grammarlens}[title={What you've built}]
One more word for getting around.
\end{grammarlens}

\subsection*{Your turn}
\begin{itemize}
\raggedright
  \item \emph{Say it:} \emph{hisāb}
  \item \emph{Say it:} \emph{hisāb}, once more
  \item \emph{Say it:} \emph{dukāndār sab grāhakoṁ kā hisāb rakhtā hai.} — The shopkeeper keeps an account for every customer.
  \item \emph{Recall:} say the Hindi for credit, then the Hindi for an instalment, then say \emph{hisāb} again
\end{itemize}

\begin{tcolorbox}[breakable,colback=teal!4,colframe=teal!35!black,title={Before you move on}]
What does \hi{हिसाब} mean? (\textquotedblleft{}An account; a reckoning\textquotedblright{}.) Say it once more.
\end{tcolorbox}
\section[\texorpdfstring{\hi{दलाल} (dalāl)}{दलाल (dalāl)}]{\hi{दलाल} (dalāl) — a broker, a middleman}

\label{lesson:HI-C427-dalal}



\begin{quote}
Before the new one: say the Hindi for an instalment, then the Hindi for an account.
\end{quote}

\subsection*{You'll want to know: \hi{दलाल}}
\textbf{\hi{दलाल}} — \emph{dalāl} — \textquotedblleft{}a broker, a middleman\textquotedblright{}.

\begin{quote}
\hi{दलाल} \hi{ने} \hi{हमें} \hi{एक} \hi{सस्ता} \hi{मकान} \hi{दिखाया।}
\end{quote}

\begin{quote}
\emph{dalāl ne hamẽ ek sastā makān dikhāyā.}
\end{quote}

\begin{quote}
The broker showed us a cheap house.
\end{quote}

Masculine or feminine, by the person meant.

\begin{cousinweb}
\textbf{\hi{दलाल}} (\emph{dalāl}) is usually traced to Arabic \emph{dallāl}, \textquotedblleft{}a broker, one who shows the way to a sale\textquotedblright{}, through Persian.
\end{cousinweb}

\begin{grammarlens}[title={What you've built}]
One more word for getting around.
\end{grammarlens}

\subsection*{Your turn}
\begin{itemize}
\raggedright
  \item \emph{Say it:} \emph{dalāl}
  \item \emph{Say it:} \emph{dalāl}, once more
  \item \emph{Say it:} \emph{dalāl ne hamẽ ek sastā makān dikhāyā.} — The broker showed us a cheap house.
  \item \emph{Recall:} say the Hindi for an instalment, then the Hindi for an account, then say \emph{dalāl} again
\end{itemize}

\begin{tcolorbox}[breakable,colback=teal!4,colframe=teal!35!black,title={Before you move on}]
What does \hi{दलाल} mean? (\textquotedblleft{}A broker, a middleman\textquotedblright{}.) Say it once more.
\end{tcolorbox}
\section[\texorpdfstring{\hi{चुकाना} (cukānā)}{चुकाना (cukānā)}]{\hi{चुकाना} (cukānā) — to pay off; to settle (a debt)}

\label{lesson:HI-C427-cukana}



\begin{quote}
Before the new one: say the Hindi for an account, then the Hindi for a broker.
\end{quote}

\subsection*{You'll want to know: \hi{चुकाना}}
\textbf{\hi{चुकाना}} — \emph{cukānā} — \textquotedblleft{}to pay off; to settle (a debt)\textquotedblright{}.

\begin{quote}
\hi{मैंने} \hi{पूरा} \hi{क़र्ज़} \hi{चुका} \hi{दिया।}
\end{quote}

\begin{quote}
\emph{maiṁne pūrā qarz cukā diyā.}
\end{quote}

\begin{quote}
I paid off the whole loan.
\end{quote}

Past, with \hi{ने} (he or she …): \hi{उसने} \hi{दुकान} \hi{का} \hi{बिल} \hi{चुकाया}, \emph{usne dukān kā bil cukāyā}.

\begin{cousinweb}
\textbf{\hi{चुकाना}} (\emph{cukānā}) goes back to \emph{*cukk}, \textquotedblleft{}fall short of, stop\textquotedblright{}, a form CDIAL reconstructs (CDIAL 4848), which lists Hindi \emph{cukānā}, \textquotedblleft{}to settle\textquotedblright{}: to bring a debt to its end.
\end{cousinweb}

\begin{grammarlens}[title={What you've built}]
One more word for getting around.
\end{grammarlens}

\subsection*{Your turn}
\begin{itemize}
\raggedright
  \item \emph{Say it:} \emph{cukānā}
  \item \emph{Say it:} \emph{cukānā}, once more
  \item \emph{Say it:} \emph{maiṁne pūrā qarz cukā diyā.} — I paid off the whole loan.
  \item \emph{Recall:} say the Hindi for an account, then the Hindi for a broker, then say \emph{cukānā} again
\end{itemize}

\begin{tcolorbox}[breakable,colback=teal!4,colframe=teal!35!black,title={Before you move on}]
What does \hi{चुकाना} mean? (\textquotedblleft{}To pay off; to settle (a debt)\textquotedblright{}.) Say it once more.
\end{tcolorbox}
